% Kamp III.2, Lemma 8: completed induction foundations and boundary decomposition.

\subsubsection{Finite families and the eleven interval boundary configurations}
\label{sec:p09-kamp-iii2-lemma8-foundations}

Kamp's Lemma~8 treats a universal disjunction whose members describe
the current point or a basic pattern connecting it to either endpoint.
Its proof uses a comparison of finite families, separates eleven boundary
configurations of a distinguished pattern, and then applies the preceding
cut construction to obtain elementary interval representations. The results
below supply the family comparison and a complete boundary decomposition
with its bracketing conditions retained
\cite[ch.~III, lemma~8, pp.~83--92]{Kamp1968}.

\begin{hbdefinition}{def:p09-kamp-universal-family}{The exact finite-family problem}
Let \(x,y,z\) be distinct variable names, and let
\[
 F\subseteq\operatorname{At}(z)\cup
                   \operatorname{Bp}(x,z)\cup\operatorname{Bp}(z,y)
\]
be finite. The unstarred \(\operatorname{At}\) is the literal class
of \cref{def:p09-kamp-iii1-components}; it contains no order atoms.
Define
\[
 D_F(z)\defeq\bigvee_{\psi\in F}\psi,\qquad
 \mathcal U_F(x,y)\defeq\forall z\,(x<z<y\to D_F(z)).
\]
An enumeration fixes the association of a nonempty disjunction; for an
empty family use \(D_\varnothing(z)=F^K_z\). Every occurrence
of \(x,y\) inside a member retains its original free-endpoint role.

For \(\psi\in\operatorname{At}\), put \(\ell(\psi)=1\).
For a basic pattern, \(\ell(\psi)\) is its occurrence-word length
from \cref{def:p09-kamp-cut-standing}. In particular a basic pattern
has odd length at least three. For a nonempty family set
\[
 M(F)\defeq\max_{\psi\in F}\ell(\psi),\qquad
 c(F)\defeq\bigl|\{\psi\in F:\ell(\psi)=M(F)\}\bigr|.
\]
For the empty family set \(M(F)=c(F)=0\).
The bars denote finite cardinality.

The source comparison is
\[
 G\prec F\quad\Longleftrightarrow\quad
 \exists m\in\mathbb N\left(
   (\forall\theta\in G\;\ell(\theta)\leq m)\ \land\
   |\{\psi\in F:\ell(\psi)=m\}|>
   |\{\theta\in G:\ell(\theta)=m\}|
                        \right).
\]
It compares finite formula sets by their lengths. It is not semantic
implication and does not count multiple occurrences of an identical
formula as distinct set members
\cite[ch.~III, p.~83]{Kamp1968}.
\end{hbdefinition}

\begin{hblemma}{thm:p09-kamp-family-rank}{proved=annex:p9.kamp.family.rank}{The source comparison and strict decrease}
For finite families in the declared classes,
\[
 G\prec F\quad\Longleftrightarrow\quad
 M(G)<M(F)\ \lor\
       \bigl(M(G)=M(F)\land c(G)<c(F)\bigr).
\]
Thus \(\prec\) is well founded. If \(\phi\) has maximal
length in a nonempty \(F\), removing \(\phi\) strictly
decreases this comparison.

More generally, remove such a \(\phi\), and replace each remaining
basic member \(\rho\) by any finite set of patterns of length at
most \(\ell(\rho)\), at most one of which has length
\(\ell(\rho)\). Retain the other length-one members, and
optionally add finitely many length-one members. If \(M(F)>1\),
every resulting family \(G\) satisfies \(G\prec F\).
This includes empty replacement sets and coincident output formulas.
The numerical condition is exactly the residual restriction proved for
both orientations of Lemma~7.
\end{hblemma}

\begin{hbdefinition}{def:p09-kamp-family-cut-data}{All choices for a finite family crossing one cut}
Let \(\rho_1,\ldots,\rho_k\), with \(k\geq0\), be
a finite list of basic patterns, with \(\rho_j\) having
\(n_j\geq1\) intervals. For each word separately, use the
cut index set \(\mathcal C_{n_j}\) and cut formulas
\(L_{j,c},R_{j,c}\) of \cref{def:p09-kamp-cut-words}.
Define the finite set of choices
\[
 \mathcal J\defeq\prod_{j=1}^k\mathcal P(\mathcal C_{n_j}).
\]
An element \(\mathbf J=(J_1,\ldots,J_k)\) chooses one
subset for each original word. For \(k=0\), \(\mathcal J\)
has one element, the empty tuple. For positive \(k\), its
cardinality is \(2^{\sum_{j=1}^k(2n_j-1)}\).

For the left-interior orientation \(\rho_j(v,y)\), define
\[
\begin{aligned}
 S_{\mathbf J}^L(u,y)&\defeq
                  \bigwedge_{j=1}^k\ \bigwedge_{c\in J_j}R_{j,c}(u,y),\\
 R_{\mathbf J}^L(v,u)&\defeq
                  \bigvee_{j=1}^k\ \bigvee_{c\in J_j}L_{j,c}(v,u).
\end{aligned}
\]
For the right-interior orientation \(\rho_j(x,v)\), define
\[
\begin{aligned}
 S_{\mathbf J}^R(x,u)&\defeq
                  \bigwedge_{j=1}^k\ \bigwedge_{c\in J_j}L_{j,c}(x,u),\\
 R_{\mathbf J}^R(u,v)&\defeq
                  \bigvee_{j=1}^k\ \bigvee_{c\in J_j}R_{j,c}(u,v).
\end{aligned}
\]
An empty total conjunction is \(T^K_u\) and an empty total
disjunction is \(F^K_v\). The two levels of indices retain
the originating pattern for every residual. \(S_{\mathbf J}\)
and \(R_{\mathbf J}\) are aggregate formula abbreviations;
they are distinct from the individual cut formulas and from the
truth set \(R_{a,b}\) introduced below.
\end{hbdefinition}

\begin{hblemma}{thm:p09-kamp-family-cut}{proved=annex:p9.kamp.family.cut}{The complete finite-family cut reduction}
Let \(x,y,u,v\) be distinct variable names and let
\(\psi_0,\psi_1\) be arbitrary formulas, with all new
internal bound names made fresh. For the left-interior patterns,
the following two formulas are equivalent over \(\mathsf{Lin}\):
\[
 \exists u\left(x<u<y\land\psi_1\land
   \forall v\left(x<v<u\to\psi_0\lor\bigvee_{j=1}^k\rho_j(v,y)\right)\right)
\]
and
\[
 \bigvee_{\mathbf J\in\mathcal J}
 \exists u\left(x<u<y\land\psi_1\land S_{\mathbf J}^L(u,y)\land
       \forall v\,(x<v<u\to\psi_0\lor R_{\mathbf J}^L(v,u))\right).
\]
For the right-interior patterns the corresponding equivalence is
between
\[
 \exists u\left(x<u<y\land\psi_1\land
   \forall v\left(u<v<y\to\psi_0\lor\bigvee_{j=1}^k\rho_j(x,v)\right)\right)
\]
and
\[
 \bigvee_{\mathbf J\in\mathcal J}
 \exists u\left(x<u<y\land\psi_1\land S_{\mathbf J}^R(x,u)\land
       \forall v\,(u<v<y\to\psi_0\lor R_{\mathbf J}^R(u,v))\right).
\]
Both statements include \(k=0\), empty individual choices,
and a varying cut at different universal witnesses. Each original
\(\rho_j\) contributes no longer residual and at most one
residual of its original length. Thus the comparison decrease of
\cref{thm:p09-kamp-family-rank} applies when the distinguished
maximal member has already been removed and the other retained
members satisfy its stated length conditions. For that application the
list enumerates each replaced family member once, and the actual
members represented by \(\psi_0\) must satisfy the rank lemma's
retention conditions. The formula equivalences themselves permit
arbitrary \(\psi_0\). Every fixed-side
conjunct remains in the output; the rank statement concerns the
residual family, not those additional conjuncts.

The method enumerates the displayed finite product and constructs
every aggregate. It performs exactly one word-cut construction per
original crossing pattern, providing the finite-family step used
when source pp.~90--91 repeat Lemma~7.
\end{hblemma}

\begin{hblemma}{thm:p09-kamp-family-normalization}{proved=annex:p9.kamp.family.normalization}{Atomic base, fresh occurrences and a removable disjunct}
If \(F\subseteq\operatorname{At}(z)\), including \(F=\varnothing\),
then \(\mathcal U_F\) has property \(I(x,y)\). An actual
elementary representative on ordered endpoints is the one-interval word
\[
                         (T^K_w,D_F(w),T^K_w).
\]

For an arbitrary finite \(F\), one can construct a family \(G\)
of the same size with the same length counts, in which every predicate
letter occurs at most once across the entire family, together with a
finite proper substitution sequence taking \(\mathcal U_G\)
back to \(\mathcal U_F\). Each \(\operatorname{At}\) member
is first represented by one fresh unary atom; each occurrence in a
basic pattern receives its own fresh letter. The substitutions restore
the complete original \(\operatorname{At}\) formulas and original
basic-pattern letters. Consequently an \(I(x,y)\) representation
for \(\mathcal U_G\) yields one for \(\mathcal U_F\).

If \(M(G)>1\), choose a further fresh letter \(Q\) and put
\(G^Q=G\cup\{Q(z)\}\). Then \(M(G^Q)=M(G)\), \(c(G^Q)=c(G)\),
and the two families have exactly the same \(\prec\)-predecessors.
Properly substituting \(F^K_w\) for \(Q\) in
\(\mathcal U_{G^Q}\) recovers a formula equivalent to
\(\mathcal U_G\). Hence an \(I(x,y)\) representation for
\(\mathcal U_{G^Q}\) also gives one for \(\mathcal U_G\).
This is the corrected falsehood substitution for the fresh-disjunct
step on source p.~84.
\end{hblemma}

\begin{hbdefinition}{def:p09-kamp-fixed-endpoint-limits}{Accumulation next to the fixed endpoint}
For a word \(W=(A_0,B_1,A_1,\ldots,B_n,A_n)\), define
\[
\begin{aligned}
 C_W^+(x)&\defeq\forall v\bigl(x<v\to
       \exists z\,(x<z<v\land\mathcal B_p(W;x,z))\bigr),\\
 C_W^-(y)&\defeq\forall v\bigl(v<y\to
       \exists z\,(v<z<y\land\mathcal B_p(W;z,y))\bigr).
\end{aligned}
\]
The first says that realizations beginning at \(x\) end
arbitrarily close to its right; the second says that realizations
ending at \(y\) begin arbitrarily close to its left.
They retain the fixed endpoint in the basic pattern.

Define labels \(\lambda_W^+(w),\lambda_W^-(w)\) as follows.
If \(n=1\), use respectively
\(\partial^+A_1(w)\) and \(\partial^-A_0(w)\).
If \(n>1\), use the actual \(\operatorname{Lf}_{n-1}(w)\)
formula with suffix word \((A_1,B_2,A_2,\ldots,B_n,A_n)\)
and the actual \(\operatorname{Lp}_{n-1}(w)\) formula with
prefix word \((A_0,B_1,A_1,\ldots,B_{n-1},A_{n-1})\),
as defined in \cref{def:p09-kamp-limiting-classes}.
Both labels belong to \(\operatorname{E}(w)\). Their
positive-stage versions require the whole indicated suffix or prefix
strictly inside every corresponding interval adjoining \(w\).
\end{hbdefinition}

\begin{hblemma}{thm:p09-kamp-fixed-endpoint-limits}{proved=annex:p9.kamp.fixed.endpoint.limits}{Elementary representations for accumulation at the fixed endpoint}
For every word \(W\) with \(n\geq1\), both \(C_W^+(x)\)
and \(C_W^-(y)\) have property \(I(x,y)\).
Explicit elementary representatives, with all labels written at their
common free argument \(w\), are
\[
\begin{aligned}
 V_W^+&=(A_0\land\lambda_W^+,\ B_1,\ T^K_w,\ T^K_w,\ T^K_w),\\
 V_W^-&=(T^K_w,\ T^K_w,\ T^K_w,\ B_n,\ A_n\land\lambda_W^-).
\end{aligned}
\]
More precisely,
\[
\begin{aligned}
 \mathsf{Lin}\models x<y\to
    (C_W^+(x)&\longleftrightarrow\mathcal B_p(V_W^+;x,y)),\\
 \mathsf{Lin}\models x<y\to
    (C_W^-(y)&\longleftrightarrow\mathcal B_p(V_W^-;x,y)).
\end{aligned}
\]
Each representative has exactly two interval pieces. The distinction
between a one-letter limiting label at \(n=1\) and a positive-stage
limiting word at \(n>1\) is retained. The first construction
supplies the source's \(\alpha_1\) step on pp.~89--90;
the second is its explicit opposite-endpoint version. Neither is a
claim about accumulation at the other endpoint while keeping a remote
initial point fixed.
\end{hblemma}

\begin{hbdefinition}{def:p09-kamp-boundary-set}{The distinguished truth set and its two orientations}
Choose a maximal-length basic member \(\phi\) of \(F\), and
put \(D(z)=D_{F\setminus\{\phi\}}(z)\). There are two
orientations. With ordered domain endpoints \(a<b\), define
\[
\begin{array}{c|c|c}
 \text{member}&H(t)&P(t)\\ \hline
 \phi(x,z)&\phi(a,t)&\exists r<t\;\phi(r,t)\\
 \phi(z,y)&\phi(t,b)&\exists r>t\;\phi(t,r).
\end{array}
\]
The table uses semantic endpoint assignment; the existential forms are
the guarded past or future closures of the same fixed word. Their
object-language versions use fresh bound endpoint names. They do not
constrain the new endpoint by \(a\) or \(b\) beyond the displayed
strict guard. Let
\[
                   R_{a,b}\defeq\{t:a<t<b\text{ and }H(t)\}.
\]
The symbol \(R_{a,b}\) is a set of domain points; it is distinct
from the cut formulas \(R_c\) of Lemma~7 and the class
\(\operatorname{E}\).
Point-argument displays below are satisfaction schemata. Their formula
versions use the original free endpoint variables in place of \(a,b\)
and fresh variables for the displayed quantified points; they add no
constants or relation symbols to \(L\). A validity display with
\(a,b\) means that its schema holds for every endpoint assignment.

For the following formulas, \(D(t)\) means evaluation at the
current-point variable \(z=t\), retaining both original endpoints
and all their occurrences inside \(D\). Write
\[
 \mathcal V_D(r,s)\defeq\forall t\,(r<t<s\to D(t)),\qquad
 \mathcal V_{D\lor P}(r,s)\defeq
                         \forall t\,(r<t<s\to D(t)\lor P(t)).
\]
These are formula abbreviations with fresh quantified variables, or
their semantic evaluations. Changing \(r,s\) changes only the
outer interval guard: it does not replace the original free \(x,y\)
inside \(D\) or the fixed endpoint inside \(H\).
\end{hbdefinition}

\begin{hbdefinition}{def:p09-kamp-boundary-types}{Exact lower and upper boundary conditions}
Assume \(R_{a,b}\ne\varnothing\). A lower boundary \(l\)
has one of the following three types:
\[
\begin{array}{c|l}
 0&l=a\text{ and }a=\inf R_{a,b},\\
 1&a<l<b\text{ and }l=\min R_{a,b},\\
 2&a<l<b,\quad l=\inf R_{a,b},\quad l\notin R_{a,b}.
\end{array}
\]
An upper boundary \(h\) has the corresponding types
\[
\begin{array}{c|l}
 0&h=b\text{ and }b=\sup R_{a,b},\\
 1&a<h<b\text{ and }h=\max R_{a,b},\\
 2&a<h<b,\quad h=\sup R_{a,b},\quad h\notin R_{a,b}.
\end{array}
\]
Infimum and supremum here are order-theoretic extrema, including
whether the boundary belongs to the set.

These are first-order conditions in the present setting. For example,
type 2 at the lower boundary is expressed by
\[
\begin{split}
 a<l<b\ \land\ \neg H(l)
 &\land\forall t\,(a<t<l\to\neg H(t))\\
 &\land\forall s\bigl(l<s\to
        \exists t\,(l<t<b\land t<s\land H(t))\bigr).
\end{split}
\]
The final quantifier ranges over all \(s>l\), so it includes
\(s=b\) and supplies an actual member. Type 2 at the upper
boundary is expressed by
\[
\begin{split}
 a<h<b\ \land\ \neg H(h)
 &\land\forall t\,(h<t<b\to\neg H(t))\\
 &\land\forall s\bigl(s<h\to
        \exists t\,(a<t<h\land s<t\land H(t))\bigr).
\end{split}
\]
Type 0 uses the corresponding coinitial or cofinal clause at \(a\)
or \(b\), with no point-membership test there. Type 1 uses
\(H(l)\) and absence of earlier members, or \(H(h)\) and
absence of later members. The annex proves every equivalence and makes
the type-0 and type-1 formulas explicit. The original fixed endpoint
in \(H\) remains fixed throughout.
\end{hbdefinition}

\begin{hblemma}{thm:p09-kamp-eleven-boundaries}{proved=annex:p9.kamp.eleven.boundaries}{The eleven complete-order configurations}
On every order in \(\mathsf{Com}\), for every \(a<b\), exactly
one of the following configurations holds:
\[
\begin{array}{c|c|c|l}
 \text{source case}&\text{lower type}&\text{upper type}&\text{additional condition}\\ \hline
 1&\multicolumn{2}{c|}{\text{empty}}&R_{a,b}=\varnothing\\
 2&0&0&l=a<h=b\\
 3&0&1&l=a<h<b\\
 4&0&2&l=a<h<b\\
 5&\multicolumn{2}{c|}{\text{singleton}}&R_{a,b}=\{c\},\ a<c<b\\
 6&1&0&a<l<h=b\\
 7&1&1&a<l<h<b\\
 8&1&2&a<l<h<b\\
 9&2&0&a<l<h=b\\
 10&2&1&a<l<h<b\\
 11&2&2&a<l<h<b.
\end{array}
\]
The bounds in each nonempty case are unique. The nine two-boundary
cases have \(l<h\); the coincident minimum and maximum belong to
the separate singleton case. These are the eleven configurations
organized by the source's \(\gamma_1,\ldots,\gamma_{11}\),
with all bounds interpreted using the same original \(H\).
The exhaustiveness assertion requires \(\mathsf{Com}\).
\end{hblemma}

\begin{hblemma}{thm:p09-kamp-hull-interpolation}{proved=annex:p9.kamp.hull.interpolation}{The fixed endpoint inside its truth-set bounds}
In either orientation, suppose the nonempty set \(R_{a,b}\)
has infimum \(l\), supremum \(h\), and \(l<h\).
Then for every \(l<t<h\),
\[
                            H(t)\longleftrightarrow P(t).
\]
There is a member of \(R_{a,b}\) strictly before \(t\) and
one strictly after \(t\); these are precisely the bracketing
premises needed for the fixed-endpoint splice. The assertion holds
on any linear order where the stated bounds exist. It does not
infer existence of those bounds on all of \(\mathsf{Lin}\).
\end{hblemma}

\begin{hbdefinition}{def:p09-kamp-corrected-boundary-branches}{The complete guarded branch formulas}
Let \(\Gamma_1\) assert emptiness of \(R_{a,b}\), and let
\(\Gamma_5(c)\) assert \(R_{a,b}=\{c\}\), with \(a<c<b\).
For each of the other nine rows let
\(\Gamma_i(l,h)\) assert exactly its lower and upper types and
its displayed strict inequalities. These symbols abbreviate the
first-order conditions just defined; they introduce no new predicates.

Define
\[
\begin{aligned}
 Z_1&\defeq\Gamma_1\land\mathcal V_D(a,b),\\
 Z_5&\defeq\exists c\bigl(\Gamma_5(c)\land
                     \mathcal V_D(a,c)\land\mathcal V_D(c,b)\bigr).
\end{aligned}
\]
For the other rows, with lower type \(\lambda\) and upper type
\(\mu\), define the endpoint tests, for \(j\in\{0,1,2\}\), by
\[
 E_j(s)\defeq
 \begin{cases}
 T^K_s,&j=0\text{ or }j=1,\\
 D(s),&j=2.
\end{cases}
\]
Next define the formula abbreviation
\[
\begin{split}
 K_i(l,h)\defeq{}&\Gamma_i(l,h)\land\mathcal V_D(a,l)
       \land E_\lambda(l)\\
 &{}\land\mathcal V_{D\lor P}(l,h)
       \land E_\mu(h)\land\mathcal V_D(h,b),
\end{split}
\]
and set
\[
 Z_i\defeq
 \begin{cases}
 K_i(a,b),&\lambda=0,\ \mu=0,\\
 \exists h\,K_i(a,h),&\lambda=0,\ \mu\ne0,\\
 \exists l\,K_i(l,b),&\lambda\ne0,\ \mu=0,\\
 \exists l\exists h\,K_i(l,h),&\lambda\ne0,\ \mu\ne0.
 \end{cases}
\]
The arguments of \(K_i\) replace only the displayed boundary
parameters \(l,h\); the original endpoints inside its constituents
stay fixed. Thus a type-0 boundary is its outer endpoint and is not
quantified. All new bound names are fresh from the original formulas.

Thus an attained interior boundary is already covered by \(H\),
an unattained interior boundary must satisfy \(D\), and an outer
endpoint is outside the original universal test. Both side intervals
require \(D\). All bracketing conditions remain in \(\Gamma_i\).
\end{hbdefinition}

\begin{hblemma}{thm:p09-kamp-corrected-boundary-decomposition}{proved=annex:p9.kamp.corrected.boundary.decomposition}{The full eleven-branch universal decomposition}
For either orientation of the distinguished maximal basic member,
\[
 \mathsf{Com}\models a<b\longrightarrow
 \left(\forall t\,(a<t<b\to D(t)\lor H(t))
                 \longleftrightarrow\bigvee_{i=1}^{11}Z_i\right).
\]
Each branch retains its actual boundary descriptor and all required
endpoint tests. Conditional on a descriptor, the corresponding branch
equivalence holds over \(\mathsf{Lin}\); complete orders are used
to guarantee that some descriptor exists. This is a corrected
boundary decomposition for the case analysis of source pp.~85--89.
\end{hblemma}

\begin{hbexample}{def:p09-kamp-bracketing-required}{A later past realization does not keep the original beginning}
Take \(T=\{0,1,2,3,4,5\}\) in its usual order, a one-interval
word \((A,B,C)\), and a remaining atom \(Q(z)\), with
\[
 A=\{0,3\},\qquad B=T\setminus\{2\},\qquad
 C=\{1,4\},\qquad Q=T\setminus\{4\}.
\]
At \(a=0,b=5\), the fixed-start truth set is \(R_{0,5}=\{1\}\).
The target universal disjunction is false at \(t=4\):
\(Q(4)\) is false and \(\phi(0,4)\) crosses the point 2
where its interval predicate fails. Nevertheless the unrestricted past
closure \(P(4)\) is true via the realization \((3,4)\).

The condition consisting only of a first witness at 1, \(Q\) below
1, and \(Q\lor P\) above 1 is therefore true. It cannot replace
the original formula. The corrected singleton branch instead requires
\(Q\) throughout \((1,5)\) and is false at 4. A branch using
\(Q\lor P\) on a two-boundary interval requires members of the
fixed-start truth set on both sides of each tested point.
\end{hbexample}

\begin{proofboundary}{def:p09-kamp-lemma8-remaining}{From boundary conditions to elementary representations}
The results above supply the comparison, normalization, guarded boundary
decomposition and fixed-endpoint accumulation formulas.
Section~\ref{sec:p09-kamp-iii2-lemma8-completion} supplies the remote
and interior approach formulas, the positive elementary supports, and the
complete smaller-family induction. Its
Theorem~\ref{thm:p09-kamp-iii2-lemma8} proves the full \(I(x,y)\)
conclusion with an explicit finite elementary disjunction.
That construction retains every fixed-side cut condition and the
bracketing needed by Lemma~6. A first-order boundary descriptor does not
acquire property \(I\) merely by being first order; the required
elementary presentations are constructed in that proof.
\end{proofboundary}
